% Fragmento público generado desde propietarios íntegros.
% Residencia: 52.4.1.
% Fuente: ../incoming/radion_monografia_20260827/manuscrito/sections/05_geometria_accion.tex:1-217
\noindent La géométrie spectrale détermine la forme et la loi surfacique détermine l'action transportée par la phase.\par\medskip

\subsubsection{Équivalence entre les cartes spectrale et géométrique}

\paragraph{Opérateur spectral positif et valeurs propres conjuguées}

Le couple angulaire produit l'opérateur bicouche
\begin{equation}
B_1=AI+C^*R,\qquad R^2=I,\qquad
P_\pm=\frac{I\pm R}{2}.
\label{eq:B1}
\end{equation}
Ses valeurs propres sont \(A\pm C^*\). Comme \(0<C^*<A\), les deux sont positives.
Dans la carte \(\kappa=\pi/180\), l'ellipse spectrale a pour demi-axes
\begin{equation}
a_1^2=\kappa(A+C^*),\qquad
b_1^2=\kappa(A-C^*).
\label{eq:semiejes-espectrales}
\end{equation}
La décomposition est réversible :
\begin{equation}
A=\frac{a_1^2+b_1^2}{2\kappa},\qquad
C^*=\frac{a_1^2-b_1^2}{2\kappa}.
\label{eq:reconstruccion-AC}
\end{equation}
L'ellipse n'est pas dessinée pour illustrer un nombre. Elle reconstruit exactement le
couple qui l'a générée.

Définissons
\begin{equation}
\rho_1=\sqrt{A^2-C^{*2}},\qquad
\eta=\operatorname{artanh}\frac{C^*}{A}
=\frac12\log\frac{A+C^*}{A-C^*}.
\label{eq:eta}
\end{equation}
Alors
\[
A\pm C^*=\rho_1e^{\pm\eta},
\]
et
\begin{equation}
a_1=\sqrt{\kappa\rho_1}\,e^{\eta/2},
\qquad
b_1=\sqrt{\kappa\rho_1}\,e^{-\eta/2}.
\label{eq:semiejes-eta}
\end{equation}
Le produit conserve l'échelle ; le quotient conserve la forme :
\[
a_1b_1=\kappa\rho_1,
\qquad
\frac{b_1}{a_1}=e^{-\eta}
=\sqrt{\frac{A-C^*}{A+C^*}}.
\]

\paragraph{Invariants focaux de l'opérateur positif}

La demi-distance focale est
\begin{equation}
c_1^2=a_1^2-b_1^2=2\kappa C^*,
\qquad
F_\pm=(\pm c_1,0).
\label{eq:focos}
\end{equation}
La distance totale vaut
\begin{equation}
d_1=2c_1=2\sqrt{2\kappa C^*}
=0.544545509325626623406299779866023\ldots.
\label{eq:d1}
\end{equation}
Le nombre isolé n'est pas la thèse. L'identité \(c_1^2=2\kappa C^*\) dit que
\(C^*\) est littéralement la scission focale quadratique dans la carte
distinguée. Les foyers sont les extrémités conjuguées de l'ouverture spectrale ;
ils ne sont pas, à eux seuls, les pôles électrique et magnétique. Cette lecture nécessite
le transducteur constitutif.

L'excentricité est
\begin{equation}
e_f=\frac{c_1}{a_1},\qquad
e_f^2=1-e^{-2\eta}=\frac{2C^*}{A+C^*}.
\label{eq:excentricidad}
\end{equation}
Numériquement,
\[
\begin{aligned}
\eta&=0.299689683921630799684926562863927\ldots,\\
e_f&=0.671451895550191986916277055864332\ldots.
\end{aligned}
\]

\paragraph{Bloc APP de référence et transport normalisé}

Le bloc central
\[
B_0=7I+2R=9P_++5P_-
\]
possède
\[
\rho_0=\sqrt{45},\qquad
\eta_0=\frac12\log\frac95,\qquad
e_0=\frac23.
\]
Sa distance focale est
\[
d_0=4\sqrt\kappa
=0.528443639680801468446003835349358\ldots.
\]
La comparaison exacte est
\begin{equation}
\boxed{\left(\frac{d_1}{d_0}\right)^2=\frac{C^*}{2}},
\qquad
d_1=d_0\sqrt{\frac{C^*}{2}}.
\label{eq:invariante-focal-relativo}
\end{equation}
Telle est la généalogie défendable de \(0.544545\ldots\) : une échelle focale APP
de référence multipliée par un invariant relatif. Sa proximité visuelle avec
\(0.54\) n'est pas une identité avec le défaut APP \(54\).

Le transport depuis \(B_0\) se sépare en échelle et anisotropie :
\begin{equation}
T_{\mathrm{rel}}=\frac{\rho_1}{\sqrt{45}}
\exp\bigl((\eta-\eta_0)R\bigr).
\label{eq:transporte-rel}
\end{equation}
Sur les feuillets,
\[
t_+=\frac{A+C^*}{9}=1.0467878786947163\ldots,
\qquad
t_-=\frac{A-C^*}{5}=1.0347228460630311\ldots.
\]
La déformation n'est pas une seule décimale. Sa partie isotrope est
\(\sqrt{t_+t_-}\) ; sa partie hyperbolique est
\(\tfrac12\log(t_+/t_-)\).

\subsubsection{Ellipse positive associée au couple angulaire}

\paragraph{Construction au moyen de demi-axes conjugués}

L'échelle spectrale précédente est dépourvue d'unités d'action. La mécanique
dimensionnelle publie la même forme sur la droite interne d'action :
\begin{equation}
a_S=\sqrt{2\hret e^\eta},
\qquad
b_S=\sqrt{2\hret e^{-\eta}}.
\label{eq:semiejes-accion}
\end{equation}
Immédiatement,
\begin{equation}
a_Sb_S=2\hret,\qquad
\frac{a_S}{b_S}=e^\eta.
\label{eq:producto-forma}
\end{equation}
Le produit fixe l'action ; le quotient fixe l'anisotropie. L'ellipse est
paramétrée par
\begin{equation}
Q(\theta)=a_S\cos\theta,\qquad
P(\theta)=b_S\sin\theta.
\label{eq:elipse-param}
\end{equation}

Les ellipses spectrale et d'action sont homothétiques :
\begin{equation}
\lambda_{\mathrm{act}}^2=\frac{2\hret}{\kappa\rho_1},
\qquad
(a_S,b_S,c_S)=\lambda_{\mathrm{act}}(a_1,b_1,c_1).
\label{eq:homotecia}
\end{equation}
Elles partagent la forme \(e^{-\eta}\), non la dimension physique.

\paragraph{Loi surfacique de l'action par radion}

La forme symplectique canonique est \(\omega=\diff Q\wedge\diff P\). L'aire
orientée balayée de \(0\) à \(\theta\) est
\begin{align}
\mathcal A(\theta)
&=\frac12\int_0^\theta
\bigl(Q(t)P'(t)-P(t)Q'(t)\bigr)\,\diff t\\
&=\frac12a_Sb_S\int_0^\theta
(\cos^2t+\sin^2t)\,\diff t.
\end{align}
En utilisant \eqref{eq:producto-forma}, nous obtenons le résultat central.

\begin{theorem}[Action par radion]
Pour l'ellipse \eqref{eq:elipse-param},
\begin{equation}
\boxed{\mathcal A(\theta)=\hret\theta.}
\label{eq:area-theorem}
\end{equation}
En particulier,
\begin{equation}
\boxed{\mathcal A(1)=\hret},
\qquad
\boxed{\mathcal A(2\pi)=2\pi\hret=\hfull}.
\label{eq:area-radion-vuelta}
\end{equation}
\end{theorem}

Cette égalité permet de définir le radion sans faire d'abord appel à la longueur
d'arc.

\begin{definition}[Radion orienté]
\begin{equation}
\mathfrak r_\sigma=(\theta=1,\sigma),
\qquad \sigma\in\{+1,-1\},
\qquad
\mathcal A(\mathfrak r_\sigma)=\sigma\hret.
\label{eq:radion-oriented}
\end{equation}
\end{definition}

Le radian est la projection qui oublie \(\sigma\), le feuillet, la mémoire et
l'aire. Le radion conserve cette information. D'où le nom \emph{rad--ion} :
l'orientation précède la charge physique, et la charge apparaît plus tard comme
\[
Q(x)=n_Q(x)e_{\mathrm{int}},\qquad
e_{\mathrm{int}}=\sqrt{\frac{2\alphah\hfull}{Z_{\mathrm{int}}}}.
\]

